% Pillar-7: controller design rules (when to escalate G, thresholds)
\section{Controller Design Rules}
\label{sec:p7-rules}
The theory and the E3 pilot jointly motivate a small set of design rules for
group-size control. We state them as rules with their supporting evidence,
in the spirit of engineering guidance rather than proven optima.

\paragraph{Rule 1: gate on PCD, alarm on ZVF.} Escalate only when ZVF is
high \emph{and} PCD indicates interior difficulty (PCD well above $0$, i.e.\
a nontrivial fraction of groups carry contrast). High ZVF with PCD
$\approx 0$ and high mean reward is mastery: escalating $G$ there buys
improvements of only $O(\eps)$ per added rollout, since
$h_G(p) \approx 1 - G\eps$ near the boundary
(Appendix~\ref{sec:p7-appendix})---rollouts spent confirming what the policy
already knows. High ZVF with PCD $\approx 0$ and near-zero mean reward is
incapacity: no feasible $G$ rescues $p \approx 0$, and the correct
intervention is the data (curriculum, decomposition), not the sampler.
The symmetry $h_G(p)=h_G(1-p)$ is why this gate is mandatory (the
cross-model rows of Table~\ref{tab:p7-ushape} illustrate both ends, though
their interior is confounded with configuration); the micro-jitter result is why the alarm must be
tie-based ZVF on the \emph{binary} sub-reward, computed before any dense
shaping terms are added.

\paragraph{Rule 2: escalate multiplicatively, in small steps.} T2 gives
exponential returns in $G$ at interior $p$ (at $p = \nicefrac12$,
$h_G = 2^{1-G}$, so each additional rollout halves the degenerate
fraction), but linear cost. This rule is a recommendation from the theory, not an E3 finding: the E3
escalator (steps of $+2$ from $G{=}4$) cost $+55\%$ rollouts and moved mean
ZVF only from $0.25$ to $0.23$ (Table~\ref{tab:p7-e3}). Jumping straight to
a large fixed $G$ pays the overhead on every step including the healthy
ones. The repo sweep
(Table~\ref{tab:p7-gsweep}) makes the same point statically: $G{=}16$ costs
$4\times$ the rollouts of $G{=}4$ for a ZVF drop from $0.76$ to $0.63$ and
no held-out gain on a near-saturated task.

\paragraph{Rule 3: trigger on spikes, not levels.} Pillar~2 reported
that ZVF is temporally sticky (lag-1 autocorrelation $\approx 0.94$) and
drifts upward over training, so a level threshold tends to fire either never
or permanently. We therefore recommend triggering on upward \emph{spikes}
relative to the run's own recent trajectory. The E3 pilot did \emph{not}
implement this: its trigger is a level threshold on a two-step moving average
($>0.4$), so E3 provides no evidence for or against spike triggering. Sticky
drift toward high ZVF with rising reward is the expected saturation endgame,
not an emergency.

\paragraph{Rule 4: never compare controllers across stacks.} The
under-specified same-label comparison ($85.6\%$ vs $5.0\%$, with backend and
base checkpoint both differing) and the DAPO label flip (mean ZVF $0.58$
closed vs $0.00$ open under
the same name) mean a controller evaluated on one stack has not been
evaluated at all until the stack is reported. Controller comparisons should
ship the seven-field run manifest of the companion papers' eight-item
minimum reporting standard---loss form,
reference-policy/KL handling, sampler/backend/precision, per-step ZVF/GU
trajectory, group-size schedule, held-out split, and
decontamination/parser probes---with the group-size \emph{schedule} (item
5) now carrying the controller's escalation log.

\paragraph{Rule 5: treat ZVF $= 0$ as a purchase, not a prize.} Dynamic
sampling reports zero ZVF among accepted groups in the E3 pilot,
at $+45\%$ rollout cost. This is not a termination or cost guarantee
on other populations. Whether to buy it is a compute-allocation
decision; in the two-seed pilot it bought no held-out gain
(Table~\ref{tab:p7-e3}). The controller framework makes the price explicit;
we expect the trade to tip toward aggressive intervention on harder prompt
populations, and flag that as the decisive open experiment
(Section~\ref{sec:p7-limitations}).

\paragraph{Endgame evidence that $G$ wants a schedule.} A two-seed
matched-budget panel (2{,}560 rollouts per arm: $G{=}2\times160$ steps
vs.\ $G{=}16\times20$ steps, Qwen3-8B/GSM8K, batch~8, seeds 123 and 456)
shows the wall the rules above exist to avoid, at the \emph{high}-accuracy
end. The $G{=}2$ arms reach last-10 train reward $0.90$ and $0.96$ and then
idle at mean ZVF $0.975$ and $0.981$ over their last 20 steps --- all-correct
groups with zero centered reward contrast in the stated runner, not
necessarily zero total gradient. The $G{=}16$ arms average ZVF $0.19$ and
$0.11$ (maxima $0.50$ and $0.38$) but, with only 20 optimizer steps, reach
last-10 train reward of just $0.32$ and $0.39$. Neither arm has a held-out
evaluation in these records. A fixed small $G$ buys more optimizer steps and
starves in the endgame; a fixed large $G$ keeps contrast but, on the same budget,
reaches far lower training reward. The
static trade is exactly what Rules~1--3 convert into a closed-loop
schedule: start small, escalate as $\hat p$ (and with it
$\mathrm{ZVF}$) rises. Artifacts:
\texttt{tinker-runs/results/er2b\_*.json}.
